% SGA 3，第 XI 篇；根据英文重建源翻译。
% 范围：命题 3.7 至推论 3.11 全部；止于命题 3.12 之前。
% 权威来源：Polo--Gille Expo11.pdf，组件局部页 13--17；合卷页 735--739；印刷／原文页边页 147--153。
% Jacob Reinhold 的 CC BY 4.0 Markdown 谱系仅用于措辞比较；文字、公式、编号、注释及图均以 Polo--Gille PDF 为准。
% 精确续接／终点：命题 3.12，Expo11 局部页 17／合卷页 739／原文页边页 153。

\medskip\noindent\SGARefTarget{sga3:xi:proposition:3:7}\textbf{命题 3.7。} --- \textit{设 \(S\) 为局部 Noether 预概形，\(I\) 为有向偏序指标集，\((T_i)_{i\in I}\) 为局部有限型 \(S\)-预概形的投射系，\[
  T=\varprojlim_{i\in I}T_i
\] 为投射极限函子，\(F\) 为 \(S\) 上的函子，\(u:F\to T\) 为 \(S\)-态射。假设下列条件成立：}

\begin{enumerate}\item[(a)]\SGARefTarget{sga3:xi:proposition:3:7:item:a}\textit{\(F\) 为忠实平坦拟紧拓扑的层，与环的归纳极限交换，并与局部 Artin 环的理想进逆极限交换。}

\item[(b)]\SGARefTarget{sga3:xi:proposition:3:7:item:b}\textit{态射 \(u:F\to T\) 为单态射。}

\item[\(b'\)]\SGARefTarget{sga3:xi:proposition:3:7:item:b-prime}\textit{态射 \(u:F\to T\) 无穷小平展。}

\item[(c)]\SGARefTarget{sga3:xi:proposition:3:7:item:c}\textit{对 \(F\) 的每个取值于域 \(k\) 的谱的点 \(\xi\)，设 \(\xi_i\in T_i(\operatorname{Spec}k)\) 为其像，\(t_i\) 为 \(T_i\) 的相应点。存在 \(i\in I\)，使得对每个 \(j\geq i\)，过渡态射 \[
  p_{i,j}:T_j\longrightarrow T_i
\] 在 \(t_j\) 处平展。}

\item[(d)]\SGARefTarget{sga3:xi:proposition:3:7:item:d}\textit{对每个在 \(S\) 上局部有限型的预概形 \(X\) 及每个 \(S\)-态射 \(X\to F\)，此态射无穷小平展的点 \(x\in X\) 的集合是开的。}\end{enumerate}

\textit{在这些条件下，\(F\) 可由 \(S\) 上局部有限型的预概形表示。}

立即指出，在 \SGARefLink{sga3:xi:section:4}{下一节} 所考虑的情形中，条件 \SGARefLink{sga3:xi:proposition:3:7:item:c}{(c)} 和 \SGARefLink{sga3:xi:proposition:3:7:item:d}{(d)} 将用 \SGARefLink{sga3:xi:corollary:3:8}{下述推论} 验证。

\medskip\noindent\SGARefTarget{sga3:xi:corollary:3:8}\textbf{推论 3.8。} --- \textit{假设 \SGARefLink{sga3:xi:proposition:3:7:item:a}{(a)}、\SGARefLink{sga3:xi:proposition:3:7:item:b}{(b)} 和 \SGARefLink{sga3:xi:proposition:3:7:item:b-prime}{\((b')\)} 成立，则下面的条件蕴含条件 \SGARefLink{sga3:xi:proposition:3:7:item:c}{(c)} 和 \SGARefLink{sga3:xi:proposition:3:7:item:d}{(d)}：}

\begin{enumerate}\item[\(c'\)]\SGARefTarget{sga3:xi:corollary:3:8:item:c-prime}\textit{各 \(T_i\) 在 \(S\) 上平滑，且过渡态射 \(p_{i,j}:T_j\to T_i\) 平滑。}

\item[\(d'\)]\SGARefTarget{sga3:xi:corollary:3:8:item:d-prime}\textit{对 \(F\) 的每个取值于域 \(k\) 的谱的点 \(\xi\)，设 \(t_i(\xi)\) 为 \(\xi\) 所定义的 \(T_i\) 的点，\(d_i(\xi)\) 为 \(T_i\) 在 \(t_i(\xi)\) 处相对于 \(S\) 的维数，并置 \[
  d(\xi)=\sup_i d_i(\xi).
\]。则：}\begin{enumerate}\item[\(1^\circ\)]\SGARefTarget{sga3:xi:corollary:3:8:item:d-prime:1}\textit{对上述每个 \(\xi\)，有 \(d(\xi)<+\infty\)；并且}

\item[\(2^\circ\)]\SGARefTarget{sga3:xi:corollary:3:8:item:d-prime:2}\textit{对每个在 \(S\) 上局部有限型的预概形 \(X\) 及每个 \(S\)-单态射 \(v:X\to F\)，\(X\) 上的函数 \[
    x\longmapsto d(\xi_x)
  \] 局部常值，其中对 \(x\in X\)，\(\xi_x\) 表示 \(v\) 诱导的 \(F\) 的取值于 \(\kappa(x)\) 的点。}\end{enumerate}\end{enumerate}

\smallskip\noindent\emph{命题 \SGARefLink{sga3:xi:proposition:3:7}{3.7} 的证明。}在条件 \SGARefLink{sga3:xi:proposition:3:7:item:c}{(c)} 下进行论证。设 \(t\) 为 \(\xi\) 定义的 \(\operatorname{ens}(F)\) 的元素（参见 \SGARefLink{sga3:xi:remark:3:6}{3.6}），置 \SGARefTarget{sga3:xi:display:3:7:local-ring} \[
  \mathscr O_t=\varinjlim_i\mathscr O_{T_i,t_i}.
\]。利用条件 \SGARefLink{sga3:xi:proposition:3:7:item:c}{(c)}，容易看出 \(\mathscr O_t\) 为 Noether 局部环（EGA \(0_{\mathrm{IV}}\)，10.3.1.3）。其剩余域 \(\kappa(t)\) 是各剩余域 \(\kappa(t_i)\) 的归纳极限，且 \(k\) 是它的扩张。若 \(\eta\) 表示 \(\operatorname{Spec}\kappa(t)\)，\(\xi\) 表示 \(\operatorname{Spec}k\)，则有交换图

\SGARefTarget{sga3:xi:diagram:3:7:residue-field-square} \[
\begin{tikzcd}[column sep=huge,row sep=large]
  \xi \arrow[r] \arrow[d] & F \arrow[d] \\
  \eta \arrow[r] & T ,
\end{tikzcd}
\]，其定义显然，留给读者。由于 \(\xi\to\eta\) 为 fpqc 拓扑的覆盖，由 \SGARefLink{sga3:xi:proposition:3:7:item:a}{(a)} 知 \(F\) 为该拓扑的层，由 \SGARefLink{sga3:xi:proposition:3:7:item:b}{(b)} 知 \(F\to T\) 为单态射，故态射 \(\eta\to T\) 唯一地通过态射 \(\eta\to F\) 分解。现在记下 \SGARefLink{sga3:xi:lemma:3:9}{下述引理}。

\medskip\noindent\SGARefTarget{sga3:xi:lemma:3:9}\textbf{引理 3.9。} --- \textit{在 \SGARefLink{sga3:xi:proposition:3:7}{3.7} 的假设下，设 \[
  T'=\operatorname{Spec}(A)
  \quad\text{且}\quad
  T'_0=\operatorname{Spec}(k(A))
\]，其中 \(T'\) 为 Noether 局部概形，\(i:T'_0\to T'\) 为典范浸入。设给定态射的交换方块：}

\SGARefTarget{sga3:xi:diagram:3:9:local-square} \[
\begin{tikzcd}[column sep=huge,row sep=large]
  T'_0 \arrow[r] \arrow[d,"i"'] & F \arrow[d] \\
  T' \arrow[r] & T .
\end{tikzcd}
\]\textit{则存在唯一的 \(T\)-态射 \(T'\to F\)。}

证明就是 \(\SGARefLink{sga3:xi:proposition:3:1}{3.1}\) 中 \((\mathrm{iii})\Rightarrow
(\mathrm{ii})\) 的证明，其中并未使用该陈述中各 \(S\) 的可表性。所用事实为：\(F\) 为 fpqc 拓扑的层，与局部 Artin 环（这里是环 \[
  A/\operatorname{rad}(A)^{n+1},
\]）的理想进逆极限交换，且 \(F\to T\) 为单态射并且无穷小平展。

在 \[
  T'=\operatorname{Spec}(\mathscr O_t),
  \qquad
  T'_0=\operatorname{Spec}(\kappa(t)).
\] 时应用引理 \SGARefLink{sga3:xi:lemma:3:9}{3.9}。我们刚刚构造了 \(T'_0\to F\)，而典范同态 \[
  \mathscr O_{T_i,t_i}\longrightarrow\mathscr O_t
\] 定义态射的投射系 \(\operatorname{Spec}(\mathscr O_t)\to T_i\)，从而定义典范态射 \(\operatorname{Spec}(\mathscr O_t)\to T\)。\SGARefTarget{sga3:xi:diagram:3:9:application-square} 相应方块 \(\mathrm{(x)}\) 显然交换：根据 \(T\) 的定义，只需在将 \(T\) 换成 \(T_i\) 后验证。因此得到唯一的 \(T\)-态射

\SGARefTarget{sga3:xi:display:3:9:v-prime} \[
  v':T'=\operatorname{Spec}(\mathscr O_t)\longrightarrow F.
\]。由于 \(F\) 与环的归纳极限交换，对充分大的 \(i\)，\SGARefLink{sga3:xi:display:3:9:v-prime}{此态射} 通过

\SGARefTarget{sga3:xi:display:3:9:v-i-prime} \[
  v'_i:T'_i=\operatorname{Spec}(\mathscr O_{T_i,t_i})
  \longrightarrow F.
\] 分解。对这样的 \(i\)，有 \SGARefTarget{sga3:xi:display:3:9:local-isomorphism} \[
  T'\xrightarrow{\sim}T'_i,
  \qquad\text{即，}\qquad
  \mathscr O_{T_i,t_i}\xrightarrow{\sim}\mathscr O_t.
\]。事实上，\(T'\to T'_i\) 忠实平坦且拟紧，因而为有效满态射，所以只需证明它为单态射。设两个态射 \(R\rightrightarrows T'\)，其中 \(R\) 可表，它们与 \(T'\to T'_i\) 的复合相同。则它们与 \(T'\to T\) 的复合也相同，而后者通过 \(T'\to T'_i\) 分解。因此，对每个 \(j\)，它们与 \(T'\to T_j\)、从而与 \(T'\to T'_j\) 的复合相同。由于 \(T'\) 为各 \(T'_j\) 在范畴 \((\mathbf{Sch})\) 中的投射极限，原来的两个态射相同。

考虑图

\SGARefTarget{sga3:xi:diagram:3:9:factorization-triangle} \[
\begin{tikzcd}[column sep=huge,row sep=large]
  T' \arrow[r,"v'"] \arrow[d] & F \arrow[d] \\
  T'_i \arrow[r] \arrow[ur,"v'_i"'] & T_i ,
\end{tikzcd}
\]，其中未标记的箭头为显然的映射。方块和上三角形交换。由于 \(T'\to T'_i\) 为满态射，下三角形也交换。

环 \(\mathscr O_{T_i,t_i}\) 是 \(T_i\) 中 \(t_i\) 的仿射开邻域的环的归纳极限。由于 \(F\) 与环的归纳极限交换，\(v'_i\) 来自定义在 \(T_i\) 中 \(t_i\) 的某个开邻域 \(U_t\) 上的态射 \SGARefTarget{sga3:xi:display:3:7:v-t} \[
  v_t:U_t\longrightarrow F
\]。必要时缩小此邻域后，\(v_t\) 为 \(T_i\)-态射。事实上，\(T_i\) 在 \(S\) 上局部有限型，因此也与环的归纳极限交换。\(v_t\) 与 \(F\to T_i\) 的复合为单态射，故 \(v_t\) 为单态射。

断言 \(v_t\) 在 \(t_i\) 处无穷小平展。由于 \(F\to T\) 无穷小平展，这等价于复合 \[
  T'_i\xrightarrow{v'_i}F\longrightarrow T
\] 在 \(t_i\) 处无穷小平展，或等价于诱导态射 \(T'_i\to T\) 无穷小平展。由于 \(T'\xrightarrow{\sim}T'_i\)，这又等价于 \(T'\to T\) 无穷小平展。这是直接的，因为此态射为无穷小平展态射 \(T'_j\to T_j\) 的投射极限。

应用尚未使用的条件 \SGARefLink{sga3:xi:proposition:3:7:item:d}{(d)}。可知 \(v_t\) 在 \(t_i\) 的某个邻域内无穷小平展。因此将 \(U_t\) 换成较小开集后，可以假设 \(v_t\) 为无穷小平展的单态射。

当 \(t\) 遍历 \(\operatorname{ens}(F)\) 时，态射族 \(v_t\) 满足 \SGARefLink{sga3:xi:proposition:3:5}{命题 3.5}，从而得到命题 \SGARefLink{sga3:xi:proposition:3:7}{3.7} 的结论。

\smallskip\noindent\emph{推论 \SGARefLink{sga3:xi:corollary:3:8}{3.8} 的证明。}假设条件 \SGARefLink{sga3:xi:corollary:3:8:item:c-prime}{\((c')\)} 和 \SGARefLink{sga3:xi:corollary:3:8:item:d-prime}{\((d')\)} 成立。用 \SGARefLink{sga3:xi:corollary:3:8:item:d-prime}{\((d')\)} 的记号，\(d_i(\xi)\) 对所有充分大的 \(i\) 为常值。因此 \(T_j\) 在 \(t_j\) 处相对于 \(T_i\) 的维数，即 \[
  d_j(\xi)-d_i(\xi),
\]，为零。所以 \(T_j\to T_i\) 在 \(t_j\) 处平展，这就证明了条件 \SGARefLink{sga3:xi:proposition:3:7:item:c}{(c)}。

上述证明现对每个 \(t\in\operatorname{ens}(F)\) 给出指标 \(i\)、\(t_i\) 的开邻域 \(U_t\) 以及态射 \(U_t\to F\)，它为单态射且在 \(t_i\) 处无穷小平展。还须说明 \SGARefLink{sga3:xi:display:3:7:v-t}{此态射} 在 \(t_i\) 的整个某个邻域内无穷小平展。在条件 \SGARefLink{sga3:xi:corollary:3:8:item:d-prime}{\((d')\)}、\SGARefLink{sga3:xi:corollary:3:8:item:d-prime:2}{\(2^\circ\)} 中取 \(X=U_t\)。必要时将 \(U_t\) 换成 \(t_i\) 在 \(U_t\) 中的连通分支，可以假设

\SGARefTarget{sga3:xi:display:3:8:dimension-constant} \[
  d(\xi_x)=d(\xi_{t_i})
  \qquad\text{对每个 }x\in U_t.
\]。而 \(d(\xi_{t_i})\) 就是 \(T_i\) 在 \(t_i\) 处相对于 \(S\) 的维数。由于此相对维数在 \(U_t\) 上恒定，对每个 \(x\in U_t\) 还有

\SGARefTarget{sga3:xi:display:3:8:relative-dimension-star} \[
  \tag{*}
  d(\xi_x)=\text{\(U_t\) 在 \(x\) 处相对于 \(S\) 的维数}.
\]。

另一方面，用 \SGARefLink{sga3:xi:proposition:3:7}{3.7} 的证明中的记号，考虑任意 Noether 局部预概形 \[
  R'=\operatorname{Spec}(A),
  \qquad
  R'_0=\operatorname{Spec}(k(A)).
\]。每个诱导的点 \(R'_0\to F\) 表示 \(t\in\operatorname{ens}(F)\)，即将 \(R'\) 的闭点映到 \(t\) 的态射 \(R'\to F\)，都唯一地通过 \(T'\) 分解。证明就是 \SGARefLink{sga3:xi:proposition:3:1}{3.1} 或引理 \SGARefLink{sga3:xi:lemma:3:9}{3.9} 的证明：使用 \(T'\to F\) 为在 \(t\) 处无穷小平展的单态射，以及 \(T'\) 为与理想进逆极限交换的 fpqc 层。

将此结果应用于 \(v_t\) 诱导的态射 \[
  \operatorname{Spec}(\mathscr O_{U_t,x})=R'\longrightarrow F
\]，以及 \(R'\) 的闭点的像 \(t_x\in\operatorname{ens}(F)\)，等价地，\(x\) 在 \(v_t\) 下的像。得到分解

\SGARefTarget{sga3:xi:display:3:8:local-factorization} \[
  \operatorname{Spec}(\mathscr O_{U_t,x})
  \longrightarrow
  \operatorname{Spec}(\mathscr O_{t_x})
  \longrightarrow F.
\]。第二个箭头在 \(t_x\) 处无穷小平展。因此，为证明复合在 \(x\) 处无穷小平展，只需证明第一个箭头在 \(x\) 处也如此。

由上面的 \SGARefLink{sga3:xi:diagram:3:9:application-square}{\(\mathrm{(x)}\)} 和公式 \SGARefLink{sga3:xi:display:3:8:relative-dimension-star}{\((*)\)}，第一个箭头由局部 \(S\)-同态诱导，它连接平滑 \(S\)-预概形 \(U_t\) 和 \(U_{t_x}\) 在 \(x\)、\(t_x\) 处的局部化；它们在这些点处相对于 \(S\) 有相同的维数 \(d\)。由 EGA I，6.5.1，此局部同态由 \(S\)-态射

\SGARefTarget{sga3:xi:display:3:8:w} \[
  w:U\longrightarrow V,
\] 诱导，其中 \(U\) 是 \(U_t\) 中 \(x\) 的开邻域，且 \(V=U_{t_x}\)。还须证明 \(w\) 在 \(x\) 处平展。此外，\(U\) 和 \(V\) 带有到 \(F\) 的单态射，再次缩小 \(U\) 后，\(w\) 为 \(F\)-态射。故 \(w\) 为单态射。现在只需证明 \SGARefLink{sga3:xi:lemma:3:10}{下述引理}。

\medskip\noindent\SGARefTarget{sga3:xi:lemma:3:10}\textbf{引理 3.10。} --- \textit{设 \(U\)、\(V\) 为两个平滑 \(S\)-预概形，它们相对于 \(S\) 的维数均为 \(d\)，且 \(w:U\to V\) 为单态射的 \(S\)-态射。则 \(w\) 为开浸入，因而平展。}

由 SGA 1，I，5.7，可归约到 \(S\) 为域的谱的情形，并可假设此域代数闭。由 SGA 1，I，5.1，只需证明 \(w\) 平展，且只需在 \(U\) 的闭点处证明。设 \(x\) 为这样的点，置 \(y=w(x)\)。选取 \(\mathscr O_{V,y}\) 的正则参数系 \(f_1,\ldots,f_d\)。则 \[
  A=
  \mathscr O_{U,x}\big/\sum_i f_i\mathscr O_{U,x}
\] 为 \[
  k=
  \mathscr O_{V,y}\big/\sum_i f_i\mathscr O_{V,y};
\] 的平凡扩张，即 \(A\simeq k\)。事实上，由于 \(w\) 为单态射，由其经基变换得到的结构态射 \[
  \operatorname{Spec}(A)\longrightarrow\operatorname{Spec}(k)
\] 也为单态射。由于 \(\mathscr O_{U,x}\) 为维数 \(d\) 的正则局部环，各 \(f_i\) 构成此环的正则参数系。立即可知 \(w\) 在 \(x\) 处平展。这就证明了引理 \SGARefLink{sga3:xi:lemma:3:10}{3.10}，从而证明了推论 \SGARefLink{sga3:xi:corollary:3:8}{3.8}。

\medskip\noindent\SGARefTarget{sga3:xi:corollary:3:11}\textbf{推论 3.11。} --- \textit{在 \SGARefLink{sga3:xi:corollary:3:8}{3.8} 的假设下，对 \(F\) 的每个在 \(S\) 上分离的拟紧开子函子 \(U\)，存在 \(i\in I\)，使得对每个 \(j\geq i\)，态射 \[
  u_j|_U:U\longrightarrow T_j
\] 为开浸入。特别地，若各 \(T_i\) 在 \(S\) 上拟仿射，则 \(F\) 的每个在 \(S\) 上拟紧，即在 \(S\) 上有限型的开子函子 \(U\)，在 \(S\) 上拟仿射。}

命题 \SGARefLink{sga3:xi:proposition:3:7}{3.7} 的证明表明，对每个 \(t\in\operatorname{ens}(F)\)，存在 \(i\in I\) 使得 \(u_i:F\to T_i\) 在 \(t\) 处为局部同构，进而对每个 \(j\geq i\)，\(u_j\) 在该点处也为局部同构。由拟紧性，可以选择与 \(x\in U\) 无关的 \(i\)。还须证明，对充分大的 \(i\)，\[
  u_i|_U:U\longrightarrow T_i
\] 为单态射。

由于 \(U\to T\) 为单态射，各等价关系 \[
  U\times_{T_i}U\ \subset\ U\times_S U
\] 的交为对角。预概形 \(U\times_S U\) 为 Noether 预概形，各 \(U\times_{T_i}U\) 为闭子预概形。因此其中一个已经是对角；等价地，\(u_i|_U\) 为单态射。这就证明了推论 \SGARefLink{sga3:xi:corollary:3:11}{3.11} 的第一个断言，第二个断言立即随之成立。


